% ============================================================
%  chapters/ch1.tex — Chương 1: Giới thiệu và tổng quan
% ============================================================
\section{GIỚI THIỆU VÀ TỔNG QUAN}
\label{sec:ch1}

% ============================================================
\subsection{Bối cảnh và động lực nghiên cứu}
\label{subsec:context}
% ============================================================

Trong những thập kỷ gần đây, \textit{Internet of Things} (IoT) đã trở thành
một trong những xu hướng công nghệ phát triển nhanh nhất, kết nối hàng tỷ
thiết bị thông minh trong các lĩnh vực từ nhà thông minh, y tế, nông nghiệp
chính xác đến hệ thống công nghiệp và thành phố thông minh \cite{prajisha2025dis}.
Theo dự báo của các tổ chức nghiên cứu công nghệ, số lượng thiết bị IoT được
kết nối trên toàn cầu sẽ vượt mốc 30 tỷ thiết bị vào năm 2030, tạo nên một
hạ tầng mạng lưới phức tạp và phân tán chưa từng có trong lịch sử.

Nền tảng vật lý của hạ tầng IoT chủ yếu dựa trên các \textit{Mạng năng lượng
thấp và mất mát cao} (Low-power and Lossy Networks -- LLN). Đây là các mạng
không dây bao gồm hàng nghìn nút cảm biến bị giới hạn nghiêm ngặt về bộ nhớ,
khả năng xử lý và nguồn năng lượng pin \cite{prajisha2025dis}. Các nút trong
LLN thường được triển khai trong những môi trường khắc nghiệt nơi việc thay
thế pin là không khả thi, đòi hỏi giao thức mạng phải cực kỳ tiết kiệm tài
nguyên.

Để giải quyết bài toán định tuyến hiệu quả trong môi trường ràng buộc này,
Tổ chức Kỹ thuật Internet (IETF), thông qua nhóm làm việc ROLL
(\textit{Routing Over Low-power and Lossy Networks}), đã chuẩn hóa \textbf{RPL}
(\textit{IPv6 Routing Protocol for Low-Power and Lossy Networks}) theo
RFC~6550 vào năm 2012 \cite{winter2012rpl}. RPL nhanh chóng trở thành
giao thức định tuyến tiêu chuẩn \textit{de facto} cho IoT/LLN nhờ thiết kế
nhẹ nhàng, hỗ trợ đa dạng mô hình truyền thông (nhiều-đến-một,
một-đến-nhiều, điểm-đến-điểm) và khả năng tự tổ chức tô-pô thích ứng.

Tuy nhiên, chính sự đơn giản và phổ biến của RPL lại tạo ra bề mặt tấn công
đáng lo ngại. Các nghiên cứu gần đây đã chỉ ra rằng RPL rất dễ bị tổn thương
trước nhiều dạng tấn công khác nhau, gây ra các hậu quả nghiêm trọng như:
tăng tiêu thụ năng lượng, định tuyến gói sai hướng, tắc nghẽn mạng và cô lập
các nút hợp lệ \cite{prajisha2025dis}. Đặc biệt nguy hiểm là các cuộc tấn
công từ \textit{nút nội bộ bị xâm phạm} (compromised internal node) ---
loại tấn công mà ngay cả cơ chế mật mã tùy chọn của RPL cũng không bảo vệ
được \cite{dvir2011vera}.

Trong bối cảnh đó, nhu cầu xây dựng \textbf{Hệ thống Phát hiện Xâm nhập}
(Intrusion Detection System -- IDS) chuyên biệt cho mạng RPL ngày càng trở
nên cấp thiết, đặc biệt khi các ứng dụng IoT trong y tế, cơ sở hạ tầng
trọng yếu và quân sự yêu cầu mức độ tin cậy và an toàn cao.

% ============================================================
\subsection{Vấn đề nghiên cứu}
\label{subsec:problem}
% ============================================================

RPL xây dựng và duy trì một đồ thị acyclic có hướng tập trung vào đích
(\textit{Destination Oriented Directed Acyclic Graph} -- DODAG) thông qua
ba loại thông điệp điều khiển chính: \texttt{DIO} (DODAG Information Object),
\texttt{DIS} (DODAG Information Solicitation) và \texttt{DAO}
(Destination Advertisement Object). Thiết kế tin tưởng ngầm định vào tính
trung thực của các nút tham gia mạng, điều này hoàn toàn có thể bị khai thác
bởi các nút độc hại.

Theo phân loại trong \cite{overview_rpl_framework}, các tấn công vào RPL
thuộc ba nhóm chính:

\begin{enumerate}
    \item \textbf{Tấn công cạn kiệt tài nguyên (Resources):} Nút độc hại
          buộc các nút hợp lệ thực hiện những hành động không cần thiết,
          làm tăng tiêu thụ năng lượng, bộ nhớ và băng thông. Nhóm này
          bao gồm tấn công trực tiếp (direct) như DIS Flooding, và gián
          tiếp (indirect) như Version Number Attack.

    \item \textbf{Tấn công phá vỡ tô-pô (Topology):} Nút độc hại làm méo
          cấu trúc DODAG, khiến mạng hội tụ về trạng thái không tối ưu
          (sub-optimization) hoặc cô lập một số nút (isolation). Nhóm
          này bao gồm Decreased Rank Attack và Blackhole Attack.

    \item \textbf{Tấn công lưu lượng (Traffic):} Nút độc hại len lỏi vào
          mạng để nghe trộm (eavesdropping) hoặc giả mạo
          (misappropriation) lưu lượng dữ liệu hợp lệ.
\end{enumerate}

Đồ án tập trung vào \textbf{bốn dạng tấn công} tiêu biểu và nguy hiểm nhất,
có đầy đủ tài liệu nghiên cứu và có thể mô phỏng kiểm thử trên Cooja:

\begin{enumerate}
    \item \textbf{DIS Flooding Attack} \textit{[Resources -- Direct]}:
          Kẻ tấn công liên tục phát tán bản tin DIS vào mạng, buộc tất cả
          các nút lân cận phải reset bộ hẹn giờ Trickle và phản hồi bằng
          bản tin DIO, dẫn đến bùng nổ lưu lượng điều khiển và cạn kiệt
          năng lượng pin \cite{prajisha2025dis}.

    \item \textbf{Blackhole Attack} \textit{[Topology -- Isolation]}:
          Nút tấn công quảng bá các thông số định tuyến hấp dẫn để thu
          hút lưu lượng về phía mình, sau đó lặng lẽ hủy toàn bộ gói tin
          thay vì chuyển tiếp, gây mất gói diện rộng và có thể cô lập
          các nhánh mạng khỏi nút gốc \cite{sharma2022blackhole}.

    \item \textbf{Decreased Rank Attack} \textit{[Topology --
          Sub-optimization]}: Nút tấn công quảng bá giá trị Rank thấp hơn
          thực tế để lôi kéo các nút lân cận chọn nó làm nút cha ưu tiên,
          từ đó kiểm soát luồng dữ liệu và làm méo cấu trúc DODAG
          \cite{decrank2022}.

    \item \textbf{VeRA / DODAG Version Number Attack} \textit{[Resources
          -- Indirect]}: Kẻ tấn công tăng số phiên bản DODAG bất hợp
          pháp, buộc toàn bộ mạng tái cấu trúc định tuyến toàn cục
          (global repair), tiêu tốn lượng lớn năng lượng và băng thông
          điều khiển \cite{dvir2011vera}.
\end{enumerate}

Mặc dù đã có các nghiên cứu đề xuất giải pháp phòng chống từng dạng tấn
công riêng lẻ, việc xây dựng một IDS thống nhất có khả năng phát hiện đồng
thời nhiều kiểu tấn công trên thiết bị cực kỳ hạn chế tài nguyên vẫn là
một thách thức mở.

% ============================================================
\subsection{Mục tiêu đề tài}
\label{subsec:objectives}
% ============================================================

Đồ án hướng tới năm mục tiêu cụ thể sau:

\begin{enumerate}
    \item[\textbf{M1.}] \textbf{Nghiên cứu lý thuyết:} Nắm vững nguyên
          lý hoạt động của RPL (DODAG, Trickle Timer, Objective Function,
          DIO/DIS/DAO), phân loại tấn công và các khái niệm nền tảng của
          IDS trong môi trường IoT/LLN.

    \item[\textbf{M2.}] \textbf{Phân tích tấn công:} Hiểu sâu cơ chế,
          tác động đến hiệu năng mạng và dấu hiệu nhận biết của 4 dạng
          tấn công RPL được chọn thông qua tài liệu và mô phỏng thực nghiệm.

    \item[\textbf{M3.}] \textbf{Thiết kế IDS:} Đề xuất kiến trúc IDS phù
          hợp với ràng buộc tài nguyên của LLN, tích hợp module phát hiện
          độc lập cho từng kiểu tấn công mà không cần sửa đổi nhân RPL.

    \item[\textbf{M4.}] \textbf{Triển khai và kiểm thử:} Cài đặt IDS
          trên Contiki-NG và mô phỏng trên Cooja theo ba trạng thái:
          baseline (bình thường), đang bị tấn công và có IDS bảo vệ.

    \item[\textbf{M5.}] \textbf{Đánh giá định lượng:} So sánh các chỉ số
          PDR, End-to-End Delay, Control Overhead, Power Consumption,
          Detection Rate (TPR) và False Positive Rate (FPR) trước và sau
          khi triển khai IDS.
\end{enumerate}

% ============================================================
\subsection{Phạm vi nghiên cứu}
\label{subsec:scope}
% ============================================================

Để đảm bảo tính khả thi trong khuôn khổ đồ án môn học, phạm vi được giới
hạn cụ thể như sau:

\begin{itemize}
    \item \textbf{Giao thức:} Chỉ nghiên cứu RPL (RFC~6550). Không bao
          gồm các giao thức LLN khác như OSPF, IS-IS hay AODV.

    \item \textbf{Mô hình kẻ tấn công:} Giả định tấn công từ \textit{nút
          nội bộ bị xâm phạm} (compromised internal node). Đây là mô
          hình nguy hiểm nhất vì kẻ tấn công đã vượt qua cơ chế mật mã
          mặc định của RPL, có khả năng gửi bản tin hợp lệ về cú pháp
          nhưng mang nội dung độc hại.

    \item \textbf{Môi trường:} Mô phỏng trên Contiki-NG và Cooja
          Simulator. Không triển khai trên phần cứng thực tế.

    \item \textbf{Phương pháp IDS:} Tập trung vào anomaly-based và
          threshold-based detection. Không đi sâu vào học máy hoặc
          học sâu.

    \item \textbf{Quy mô mô phỏng:} Mạng tĩnh từ 10 đến 30 nút, tỉ lệ
          nút tấn công từ 0\% đến 50\%, thời gian mô phỏng 300 giây
          mỗi kịch bản.
\end{itemize}

% ============================================================
\subsection{Đóng góp của đề tài}
\label{subsec:contribution}
% ============================================================

So với các nghiên cứu liên quan, đồ án có bốn đóng góp chính:

\begin{enumerate}
    \item \textbf{Khung phân tích thống nhất:} Tổng hợp và phân tích chi
          tiết 4 dạng tấn công RPL trong một khung đánh giá nhất quán với
          cùng bộ chỉ số hiệu năng, cho phép so sánh trực tiếp mức độ
          nguy hại của từng kiểu tấn công.

    \item \textbf{Kiến trúc IDS tích hợp:} Đề xuất IDS có khả năng phát
          hiện đồng thời nhiều kiểu tấn công mà không yêu cầu sửa đổi
          nhân RPL, đảm bảo tương thích với các triển khai Contiki-NG
          hiện có.

    \item \textbf{Bộ kịch bản kiểm thử chuẩn hóa:} Cung cấp bộ kịch
          bản Cooja có thể tái lập (reproducible), giúp cộng đồng nghiên
          cứu dễ dàng so sánh và mở rộng.

    \item \textbf{Phân tích đánh đổi tài nguyên:} Làm rõ sự đánh đổi
          (trade-off) giữa độ chính xác phát hiện và chi phí tài nguyên
          phát sinh bởi IDS trên các nút cảm biến hạn chế.
\end{enumerate}

% ============================================================
\subsection{Cấu trúc báo cáo}
\label{subsec:structure}
% ============================================================

Báo cáo được tổ chức thành 6 chương như sau:

\begin{itemize}
    \item \textbf{Chương 1 -- Giới thiệu và tổng quan} (chương hiện tại):
          Trình bày bối cảnh, vấn đề nghiên cứu, mục tiêu, phạm vi và
          cấu trúc báo cáo.

    \item \textbf{Chương 2 -- Cơ sở lý thuyết:} Giới thiệu chi tiết giao
          thức RPL, phân loại tấn công theo taxonomy, tổng quan về IDS
          trong IoT, các chỉ số đánh giá hiệu năng và môi trường
          Contiki-NG/Cooja.

    \item \textbf{Chương 3 -- Phân tích các dạng tấn công RPL:} Mô tả
          chi tiết cơ chế, tác động và dấu hiệu nhận biết của 4 kiểu
          tấn công, kèm kết quả mô phỏng so sánh với baseline.

    \item \textbf{Chương 4 -- Thiết kế IDS phát hiện tấn công RPL:}
          Khảo sát related work, đề xuất kiến trúc IDS và trình bày
          các module phát hiện/phòng chống tương ứng.

    \item \textbf{Chương 5 -- Demo và kiểm thử:} Trình bày môi trường
          thực nghiệm, kịch bản kiểm thử, kết quả đo lường (PDR, E2ED,
          overhead, power) và đánh giá IDS (TPR, FPR).

    \item \textbf{Chương 6 -- Kết luận và hướng phát triển:} Tóm tắt
          kết quả, hạn chế và các hướng nghiên cứu tiếp theo.
\end{itemize}

Ngoài ra, báo cáo bao gồm \textbf{Phụ lục} với mã nguồn IDS
(Contiki-NG), script mô phỏng Cooja và danh mục chữ viết tắt.